\documentclass[10pt,a4paper]{amsart}
\usepackage[margin=19mm]{geometry}
\usepackage{amsmath,amssymb,mathtools}
\usepackage[scr=rsfs,cal=euler]{mathalfa}
\usepackage{xfrac}
\usepackage[unicode,hidelinks,bookmarks=false]{hyperref}
\newcommand{\ie}{i.e.\ }
\newcommand{\cf}{cf.\ }
\newcommand{\resp}{respectively\ }
\newcommand{\Subsection}{\subsection}
\providecommand{\nobreakdash}{\nobreak\hbox{-}}
\setlength{\emergencystretch}{2em}
\multlinegap0pt
\usepackage{mathtools}
\usepackage{amssymb}
\usepackage{xcolor}
\usepackage{float}
\usepackage[normalem]{ulem}
\newtheorem{lemma}{Lemma}
\newtheorem{proposition}{Proposition}
\newtheorem{definition}{Definition}
\newcommand{\PS}[1]{$(\text{PS})_{#1}$}
\def \R {\mathbb{R}}
\renewcommand{\ge}{\geqslant}
\newcommand{\half}{\frac{1}{2}}
\renewcommand{\le}{\leqslant}
\title{An existence theory \\ for superposition operators of mixed order\\ subject to jumping nonlinearities}
\author{Serena Dipierro, Kanishka Perera, Caterina Sportelli, Enrico Valdinoci}
\date{Source: arXiv:2309.13895v1 (CC BY 4.0)}
\begin{document}
\maketitle

\noindent\emph{Source and licence.} An existence theory \\ for superposition operators of mixed order\\ subject to jumping nonlinearities. Version arXiv:2309.13895v1 (CC BY 4.0), distributed by arXiv under the Creative Commons Attribution 4.0 International licence, http://creativecommons.org/licenses/by/4.0/, as recorded in the arXiv licence metadata for this exact version and shown on the abstract page https://arxiv.org/abs/2309.13895v1. Full licence notice: \texttt{licenses/arxiv-2309.13895-CC-BY-4.0.txt}.\par\smallskip

\noindent\textit{Editorial scope.} Counted unit: the article's proposition P51 (source0.tex lines 766--776). The article's complete original proof of it is delivered with it (source0.tex lines 778--975, one \texttt{proof} environment of 198 lines). No mathematics is written, repaired or re-derived: the statement and the proof are the article's own lines, in the article's own notation, and no retained line is substituted or re-flowed.  Retained uncounted as the source-local context the proof needs: lines 224--226; lines 227--229; lines 231--234; lines 243--245; lines 254--259; lines 334--340; lines 537--540; lines 613--620; lines 649--658; lines 671--680; lines 683--686; lines 702--716.  Nothing was omitted from any retained span, so the package carries no omission record.  No retained line contains a citation, so the wrapper carries no bibliography and no bibliography entry.  No retained line contains a figure environment, an image-inclusion command or drawing code, so no image is delivered and the package declares no asset dependency.  Editorial formatting only, in addition: an AMS \texttt{amsart} preamble that restates the article's own declarations and macros used by the retained lines, namely the environments \texttt{lemma}, \texttt{proposition}, \texttt{definition} on separate counters, together with the standard packages those lines need.  The retained environments print in this document as Lemma 1, Proposition 1, Definition 1 in order of appearance, whereas the article prints the same items with the numbering of its own full document.  The complete native source of the article as distributed in the arXiv e-print for this exact version is bundled unchanged as \texttt{sources/147502/source0.tex}.
\begin{equation}\label{mu00}
{\mu^+}([\overline s, 1])>0,
\end{equation}

\begin{equation}\label{mu3}
{\mu^-}_{\big|_{[\overline s, 1]}}=0
\end{equation}

\begin{equation}\label{mu2}
\mu^-\big([{{ 0 }}, \overline s]\big)\le\gamma
\mu^+\big([\overline s, 1]\big).
\end{equation}

\begin{equation}\label{scritico}
\mu^+ ([s_\sharp, 1])> 0.
\end{equation}

\begin{equation} \label{mainab}
\left\{\begin{aligned}
\int_{[{{ 0 }}, 1]} (-\Delta)^s u\, d\mu(s) & = bu^+ - au^- + |u|^{2_{s_\sharp}^\ast - 2}\, u && \text{in } \Omega,\\
u & = 0 && \text{in } \R^N \setminus \Omega,
\end{aligned}\right.
\end{equation}

\begin{equation}\label{best_rho}\begin{split}{\mathcal{S}}:=&\inf\Bigg\{
\mu^+(0)\,\|u\|^2_{L^2(\Omega)}
+\mu^+(1)\,\|\nabla u\|^2_{L^2(\Omega)}
\\&\qquad\qquad+
\int_{({{ 0 }},1)}\left[
{c_{N,s}} \iint_{\R^{2N}} \frac{| u(x)-u(y)|^2}{|x-y|^{N+2s}}\, dx\,dy\right]\,d\mu^+(s)
\Bigg\}.\end{split}\end{equation}

\begin{lemma}\label{hilbert}
Suppose that~$\mu^+$ satisfies~\eqref{mu00}.
Then, $\mathcal{X}(\Omega)$ is a Hilbert space.
\end{lemma}

\begin{proposition}\label{crucial} Assume~\eqref{mu3} and~\eqref{mu2}.

Then, there exists~$c_0=c_0(N,\Omega)>0$ such that,
for any~$u\in\mathcal{X}(\Omega)$, we have
\[
\int_{[{{ 0 }}, \overline s]} [u]_{s}^2 \, d\mu^- (s) \le c_0\,\gamma \int_{[\overline s, 1]} [u]^2_{s} \, d\mu(s).
\]
\end{proposition}

\begin{proposition}\label{emb}Assume~\eqref{mu00}, \eqref{mu3} and~\eqref{mu2}.
Let~$s_\sharp\in [\overline s, 1]$ be as in~\eqref{scritico}. 

Then, there exists a positive constant~$\bar c =\bar c (N,\Omega, s_\sharp)$ such that, for any~$u\in \mathcal{X}(\Omega)$,
\begin{equation}\label{pqwodfl134rt}
[u]_{s_\sharp}\le \bar{c}\left(\, \int_{[{{ 0 }}, 1]} [u]^2_{s}\, d\mu^+ (s)\right)^{\frac12}.
\end{equation}
In particular, the space~$\mathcal{X}(\Omega)$ is continuously embedded in~$L^r(\Omega)$ for any~$r\in [1, 2^*_{s_\sharp}]$ and
compactly embedded in~$L^r(\Omega)$ for any~$r\in [1, 2^*_{s_\sharp})$.
\end{proposition}

\begin{definition}\label{wsol2}
A weak solution of problem~\eqref{mainab} is a function~$u \in \mathcal{X}(\Omega)$ such that, for all~$ v \in \mathcal{X}(\Omega)$,
\[
\begin{split}
&\int_{[{{ 0 }}, 1]}\left(c_{N,s}\iint_{\R^{2N}} \frac{(u(x) - u(y))\, (v(x) - v(y))}{|x - y|^{N+2s}}\, dx\, dy\right) d\mu^+(s)\\
&\qquad- \int_{[{{ 0 }}, \overline s]}\left(c_{N,s}\iint_{\R^{2N}} \frac{(u(x) - u(y))\, (v(x) - v(y))}{|x - y|^{N+2s}}\, dx\, dy\right) d\mu^-(s)\\
&\qquad\qquad= \int_{\Omega} (bu^+ - a u^-) v\, dx + \int_\Omega |u|^{2_{s_\sharp}^\ast - 2}\, uv\, dx.
\end{split}
\]
\end{definition}

\begin{equation}\label{fun}
E(u) = \half \,\int_{[{{ 0 }}, 1]} [u]^2_{s}\, d\mu^+ (s)
 -\half\int_{[{{ 0 }}, \overline s]} [u]^2_s\, d\mu^-(s) - \half \int_\Omega \left[a\, (u^-)^2 + b\, (u^+)^2\right]\, dx - \frac{1}{2_{s_\sharp}^*} \int_\Omega |u|^{2_{s_\sharp}^*}\, dx.
\end{equation}

\begin{lemma}\label{Vitali}
Let~$u_n$ be a bounded sequence in~$\mathcal X(\Omega)$. 

Then, there exists~$u:\R^N\to\R$ such that, up to a subsequence, for any~$v\in\mathcal X(\Omega)$,
\begin{equation}\label{Vconv}
\begin{split}
&\lim_{n\to+\infty}\int_{[{{ 0 }}, \overline s]} \left(\;\iint_{\R^{2N}} \frac{c_{N,s}(u_n(x)-u_n(y)) (v(x)-v(y))}{|x-y|^{N+2s}} \, dx\, dy\right)\, d\mu^-(s)  \\
&\qquad =\int_{[{{ 0 }}, \overline s]} \left(\;\iint_{\R^{2N}} \frac{c_{N,s}(u(x)-u(y)) (v(x)-v(y))}{|x-y|^{N+2s}} \, dx\, dy\right)\, d\mu^-(s).
\end{split}
\end{equation}
Also, 
\begin{equation}\label{lunoconv}
{\mbox{$u_n$ converges to~$u$ in~$L^1(\Omega)$ as~$n\to+\infty$.}}
\end{equation}
\end{lemma}

\begin{proposition}\label{P51}
Let~$\mathcal S$ be as in~\eqref{best_rho} and
\begin{equation}\label{CASTA}
c^\ast:= \frac{s_\sharp}{N}\, \big((1-\theta_0)\mathcal S \big)^{\frac{N}{2 s_\sharp}}.
\end{equation}

Then, there exists~$\gamma_0>0$, depending on~$N$, $\Omega$,
$s_\sharp$, $a$ and~$b$, such that if~$\gamma\in[0,\gamma_0]$
and~$c \in (0,c^\ast)$,
then every~\PS{c} sequence of the functional~\eqref{fun} has a subsequence that converges weakly to a nontrivial critical point of~\eqref{fun}.
\end{proposition}

\begin{proof}
Let~$u_n$ be a~\PS{c} sequence of the functional~$E$, i.e.
\begin{equation}\label{c1}
\begin{split}&
\lim_{n\to+\infty}
E(u_n) \\
= \;&\lim_{n\to+\infty}\half  \int_{[{{ 0 }}, 1]} [u_n]^2_{s}\, d\mu^+ (s)
-\half\int_{[{{ 0 }}, \overline s]} [u_n]^2_s \, d\mu^-(s)
\\&\qquad - \half \int_\Omega \left[a\, (u_n^-)^2 + b\, (u_n^+)^2\right] \,dx - \frac{1}{2_{s_\sharp}^*} \int_\Omega |u_n|^{2_{s_\sharp}^*}\, dx\\= \;&c
\end{split}
\end{equation}
and~$ dE (u_n)$ converges to~$0$ in the dual of~$\mathcal{X}(\Omega)$,
namely
\begin{equation}\label{c2}
\lim_{n\to+\infty}\sup_{ v \in \mathcal{X}(\Omega)}\big|\langle dE (u_n), v\rangle\big| =0.
\end{equation}
Since, for all~$v \in \mathcal{X}(\Omega)$,
\begin{equation*}\begin{split}
\langle dE (u_n), v\rangle =
\;& \int_{[{{ 0 }}, 1]}\left(c_{N,s}
\iint_{\R^{2N}} \frac{(u_n(x) - u_n(y))\, (v(x) - v(y))}{|x - y|^{N+2s}}\, dx\, dy\right) d\mu^+(s)\\
&\qquad- \int_{[{{ 0 }}, \overline s]}\left(c_{N,s}
\iint_{\R^{2N}} \frac{(u_n(x) - u_n(y))\, (v(x) - v(y))}{|x - y|^{N+2s}}\, dx\, dy\right) d\mu^-(s)\\
&\qquad -\int_\Omega \big(b\, u_n^+ -a\, u_n^-\big)v \, dx - \int_\Omega |u_n|^{2_{s_\sharp}^\ast - 2}\, u_n v\, dx,
\end{split}
\end{equation*}
choosing~$v:=u_n$ in~\eqref{c2},
we obtain that
\begin{equation}\label{hufewghwutgui0987654}
\begin{split}0=\;& \lim_{n\to+\infty}
\langle dE (u_n), u_n\rangle\\=
\;& \lim_{n\to+\infty}\int_{[{{ 0 }}, 1]}\left(c_{N,s}
\iint_{\R^{2N}} \frac{|u_n(x) - u_n(y)|^2}{|x - y|^{N+2s}}\, dx\, dy\right) d\mu^+(s)\\
&\qquad- \int_{[{{ 0 }}, \overline s]}\left(c_{N,s}
\iint_{\R^{2N}} \frac{|u_n(x) - u_n(y)|^2}{|x - y|^{N+2s}}\, dx\, dy\right) d\mu^-(s)\\
&\qquad -\int_\Omega \big[b\, (u_n^+)^2 +a\, (u_n^-)^2\big]\, dx - \int_\Omega |u_n|^{2_{s_\sharp}^\ast}\, dx\\=\;&  \lim_{n\to+\infty}
\int_{[{{ 0 }}, 1]} [u_n]^2_{s}\, d\mu^+ (s) -\int_{[{{ 0 }}, \overline s]} [u_n]^2_s \, d\mu^-(s) - \int_\Omega \left[a\, (u_n^-)^2 + b\, (u_n^+)^2\right] \,dx - \int_\Omega |u_n|^{2_{s_\sharp}^*}\, dx
\\=\;&  \lim_{n\to+\infty} 2E(u_n)+\left(\frac2{2^*_{s_\sharp}}-1\right)\int_\Omega |u_n|^{2_{s_\sharp}^*}\, dx.
\end{split}
\end{equation}
Combining this and~\eqref{c1}, we infer that
\begin{equation}\label{jdewfguewogewuootr6u768i0}
0=2c+
\left(\frac2{2^*_{s_\sharp}}-1\right)\lim_{n\to+\infty}\int_\Omega |u_n|^{2_{s_\sharp}^*}\, dx,\end{equation}
yielding that, for large~$n$,
\begin{equation}\label{prima}
\left(\half -  \frac{1}{2_{s_\sharp}^*}\right) \|u_n\|^{2^*_{s_\sharp}}_{L^{2^*_{s_\sharp}}(\Omega)} \le c+1. 
\end{equation}
Moreover, from this and the H\"older inequality, it follows that
\begin{equation}\label{prima22}\begin{split}&
\int_\Omega \left[a\, (u_n^-)^2 + b\, (u_n^+)^2\right]\, dx \le \max\{ a, b\}\|u_n\|^2_{L^2(\Omega)} \le \max\{ a, b\} |\Omega|^{(2 s_\sharp)/N} \|u_n\|^2_{L^{2^*_{s_\sharp}}(\Omega)}\\
&\qquad\le \left(\frac{(c+1)N}{s_\sharp}\right)^{\frac2{2^*_{s_\sharp}}} \max\{ a, b\} |\Omega|^{(2 s_\sharp)/N}.
\end{split}\end{equation}

Now, by~\eqref{c1} and
Proposition~\ref{crucial}, we have that, as soon as~$n$ is big enough,
\begin{eqnarray*}
&&\frac12\int_{[{{ 0 }}, 1]} [u_n]^2_{s}\, d\mu^+ (s)\\
&\le&\half\int_{[{{ 0 }}, \overline s]} [u_n]^2_s \, d\mu^-(s) +\half \int_\Omega \left[a\, (u_n^-)^2 + b\, (u_n^+)^2\right] \,dx +\frac{1}{2_{s_\sharp}^*} \int_\Omega |u_n|^{2_{s_\sharp}^*}\, dx+c+1\\&\le&
\frac{c_0(N,\Omega)\gamma}2\int_{[ \overline s,1]} [u_n]^2_s \, d\mu(s) +\half \int_\Omega \left[a\, (u_n^-)^2 + b\, (u_n^+)^2\right] \,dx +\frac{1}{2_{s_\sharp}^*} \int_\Omega |u_n|^{2_{s_\sharp}^*}\, dx+c+1,
\end{eqnarray*}
and therefore, if~$\gamma$ is sufficiently small (possibly depending on~$N$ and~$\Omega$),
$$\frac14 \int_{[{{ 0 }}, 1]} [u_n]^2_{s}\, d\mu^+ (s)\le
\half \int_\Omega \left[a\, (u_n^-)^2 + b\, (u_n^+)^2\right] \,dx +\frac{1}{2_{s_\sharp}^*} \int_\Omega |u_n|^{2_{s_\sharp}^*}\, dx+c+1.
$$
{F}rom this, \eqref{prima} and~\eqref{prima22}, we obtain that
$$
\frac14 \int_{[{{ 0 }}, 1]} [u_n]^2_{s}\, d\mu^+ (s)\le
\half \left(\frac{(c+1)N}{s_\sharp}\right)^{\frac2{2^*_{s_\sharp}}} \max\{ a, b\} |\Omega|^{(2 s_\sharp)/N} +\frac{N(c+1)}{2s_\sharp},
$$
which says that~$ \int_{[{{ 0 }}, 1]} [u_n]^2_{s}\, d\mu^+ (s)$
is uniformly bounded in~$n$.

Hence, in view of Lemma~\ref{hilbert} and
Proposition~\ref{emb}, there exists~$u\in \mathcal{X}(\Omega)$ such that, up to subsequences,
%\begin{equation}\label{Lrs}
%u_n\to u \text{ in } L^{2}(\Omega).
%\end{equation}
\begin{equation}\label{weakun}\begin{split}
u_n\rightharpoonup u &\text{ in } \mathcal{X}(\Omega),\\ 
u_n\to u &\text{ in } L^{r}(\Omega) \text{ for every } r\in [1, 2_{s_\sharp}^*),\\ 
u_n\to u &\text{ a.e. in } \Omega.
\end{split}\end{equation}
Furthermore, we observe that~$u$ is a weak solution of~\eqref{mainab},
according to Definition~\ref{wsol2}, thanks to the convergence statements
in~\eqref{weakun} and Lemma~\ref{Vitali}.

It remains to prove that
\begin{equation}\label{uzerosper00}
u\not\equiv 0.\end{equation}
To this end,
suppose by contradiction that~$u\equiv 0$. 
We recall from~\eqref{hufewghwutgui0987654} that
\begin{eqnarray*} 0&=&
\lim_{n\to+\infty}\langle dE(u_n), u_n\rangle \\
&=&\lim_{n\to+\infty} \int_{[{{ 0 }}, 1]} [u_n]^2_{s}\, d\mu^+ (s) -\int_{[{{ 0 }}, \overline s]} [u_n]^2_s \, d\mu^-(s) - \int_\Omega \left[a\, (u_n^-)^2 + b\, (u_n^+)^2\right] \,dx - \int_\Omega |u_n|^{2_{s_\sharp}^\ast}\, dx \\&=&
\lim_{n\to+\infty} \int_{[{{ 0 }}, 1]} [u_n]^2_{s}\, d\mu^+ (s) -\int_{[{{ 0 }}, \overline s]} [u_n]^2_s \, d\mu^-(s) - \int_\Omega |u_n|^{2_{s_\sharp}^\ast}\, dx.
\end{eqnarray*}
Thus, exploiting Proposition~\ref{crucial}, we have that
\begin{eqnarray*} 0&\ge&
\lim_{n\to+\infty} \big(1-c_0(N,\Omega)\gamma\big)\int_{[{{ 0 }}, 1]} [u_n]^2_{s}\, d\mu^+ (s)  - \int_\Omega |u_n|^{2_{s_\sharp}^\ast}\, dx.
\end{eqnarray*}

Accordingly, recalling the definition of~${\mathcal{S}}$ in~\eqref{best_rho}, we infer that
\begin{eqnarray*}
0 &\ge& \lim_{n\to+\infty} 
\big(1-c_0(N,\Omega)\gamma\big)\int_{[{{ 0 }}, 1]} [u_n]^2_{s}\, d\mu^+ (s)
-\mathcal{S}^{-\frac{2_{s_\sharp}^*}2}
\left(\, \int_{[{{ 0 }}, 1]} [u_n]^2_{s}\, d\mu^+ (s)\right)^{\frac{2_{s_\sharp}^*}2}\\
&=&  \big(1-c_0(N,\Omega)\gamma\big)\\&&\qquad\times\lim_{n\to+\infty} 
\int_{[{{ 0 }}, 1]} [u_n]^2_{s}\, d\mu^+ (s) \left(1 - \frac{\mathcal{S}^{-\frac{2_{s_\sharp}^*}2}}{\big(1-c_0(N,\Omega)\gamma\big)}
\left(\, \int_{[{{ 0 }}, 1]} [u_n]^2_{s}\, d\mu^+ (s)\right)^{\frac{2_{s_\sharp}^*}2-1 }\right).
\end{eqnarray*}

Now, choosing~$\gamma$ sufficiently small (possibly in dependence of~$N$
and~$\Omega$) so that~$1-c_0(N,\Omega)\gamma>0$, we conclude that
\begin{equation}\label{jdiweogbfgsdkgjk00}
0\ge \lim_{n\to+\infty} \int_{[{{ 0 }}, 1]} [u_n]^2_{s}\, d\mu^+ (s)
\left(1 - \frac{\mathcal{S}^{-\frac{2_{s_\sharp}^*}2}}{\big(1-c_0(N,\Omega)\gamma\big)}
\left(\, \int_{[{{ 0 }}, 1]} [u_n]^2_{s}\, d\mu^+ (s)\right)^{\frac{2_{s_\sharp}^*}2-1 }\right).
\end{equation}

We observe that
\begin{equation}\label{jdiweogbfgsdkgjk}
\liminf_{n\to+\infty} \int_{[{{ 0 }}, 1]} [u_n]^2_{s}\, d\mu^+ (s)>0.
\end{equation}
Indeed, suppose by contradiction that
\begin{equation*}
\liminf_{n\to+\infty} \int_{[{{ 0 }}, 1]} [u_n]^2_{s}\, d\mu^+ (s)=0.
\end{equation*}
Then, by Proposition~\ref{crucial} we would also have that
\begin{equation*}
\liminf_{n\to+\infty} \int_{[{{ 0 }}, 1]} [u_n]^2_{s}\, d\mu^- (s)=0,
\end{equation*}
and therefore, by~\eqref{c1},
\begin{equation*}
0<c=\liminf_{n\to+\infty}\left(-\frac{1}{2_{s_\sharp}^*} \int_\Omega |u_n|^{2_{s_\sharp}^*}\, dx\right)\le 0,
\end{equation*}
which is a contradiction and thus establishes~\eqref{jdiweogbfgsdkgjk}.

Thanks to~\eqref{jdiweogbfgsdkgjk00} and~\eqref{jdiweogbfgsdkgjk},
we conclude that
\begin{equation*}
0\ge \limsup_{n\to+\infty} 
\left(1 - \frac{\mathcal{S}^{-\frac{2_{s_\sharp}^*}2}}{\big(1-c_0(N,\Omega)\gamma\big)}
\left(\, \int_{[{{ 0 }}, 1]} [u_n]^2_{s}\, d\mu^+ (s)\right)^{\frac{2_{s_\sharp}^*}2-1 }\right),
\end{equation*}
which in turn gives that
\begin{equation}\label{dhuwegi4utg4u3gtu43}
\liminf_{n\to+\infty}\int_{[{{ 0 }}, 1]} [u_n]^2_{s}\, d\mu^+ (s)\ge\big(1-c_0(N,\Omega)\gamma\big)^{\frac2{2_{s_\sharp}^*-2} }
\mathcal{S}^{\frac{2_{s_\sharp}^*}{2_{s_\sharp}^*-2}}
=\big(1-c_0(N,\Omega)\gamma\big)^{\frac{N-2s_\sharp}{2s_\sharp} }
\mathcal{S}^{\frac{N}{2_{s_\sharp}}}.
\end{equation}

Additionally, using again~\eqref{c1}, and recalling the strong convergence
statement in~\eqref{weakun},
\begin{equation*}c=
\lim_{n\to+\infty}\half  \int_{[{{ 0 }}, 1]} [u_n]^2_{s}\, d\mu^+ (s)
-\half\int_{[{{ 0 }}, \overline s]} [u_n]^2_s \, d\mu^-(s)- \frac{1}{2_{s_\sharp}^*} \int_\Omega |u_n|^{2_{s_\sharp}^*}\, dx.
\end{equation*}
Hence, exploiting Proposition~\ref{crucial}, this gives that
\begin{equation*}c\ge
\lim_{n\to+\infty}\frac{1-c_0(N,\Omega)\gamma}2
\int_{[{{ 0 }}, 1]} [u_n]^2_{s}\, d\mu^+ (s)
- \frac{1}{2_{s_\sharp}^*} \int_\Omega |u_n|^{2_{s_\sharp}^*}\, dx.
\end{equation*}
{F}rom this and~\eqref{jdewfguewogewuootr6u768i0} it follows that
\begin{equation*}c\ge
\lim_{n\to+\infty}\frac{1-c_0(N,\Omega)\gamma}2
\int_{[{{ 0 }}, 1]} [u_n]^2_{s}\, d\mu^+ (s)
- \frac{(N-2s_\sharp)c}{2{s_\sharp}},
\end{equation*}
and therefore
\begin{equation*}\frac{Nc}{2s_\sharp}\ge
\lim_{n\to+\infty}\frac{1-c_0(N,\Omega)\gamma}2
\int_{[{{ 0 }}, 1]} [u_n]^2_{s}\, d\mu^+ (s).
\end{equation*}
This and~\eqref{dhuwegi4utg4u3gtu43} give that
\begin{equation*}\frac{Nc}{s_\sharp}\ge
\big(1-c_0(N,\Omega)\gamma\big)^{\frac{N}{2s_\sharp}}
\mathcal{S}^{\frac{N}{2_{s_\sharp}}},
\end{equation*}
and thus, recalling the definition of~$c^\ast$ in~\eqref{CASTA},
\begin{equation*}c^\ast>c\ge
\big(1-c_0(N,\Omega)\gamma\big)^{\frac{N}{2s_\sharp}}\frac{s_\sharp}{N}
\mathcal{S}^{\frac{N}{2_{s_\sharp}}}
=\left(\frac{1-c_0(N,\Omega)\gamma}{1-\theta_0}\right)^{\frac{N}{2s_\sharp}} c^\ast.
\end{equation*}

We point out that this implies that
$$ c_0(N,\Omega)\gamma\ge \theta_0,$$
hence, choosing~$\gamma$ sufficiently small, possibly in dependence of~$N$, $\Omega$, $s_\sharp$, $a$ and~$b$,
we obtain the desired contradiction.

Then, the claim in~\eqref{uzerosper00} is established, completing the
proof of Proposition~\ref{P51}.
\end{proof}

\end{document}
